\chapter{Complexity and Game Metrics Analysis}
\label{game_metrics}

The previous chapter described the implementation used to simulate games of Tricolor efficiently. This chapter uses this framework to analyze the computational and gameplay properties of the game. The objective is not yet to evaluate strong AI agents, but rather to obtain a general quantitative understanding of Tricolor as a combinatorial board game.

The analysis presented here is based on large-scale simulations between Random Agents. This choice makes it possible to estimate general properties of the game without introducing strategic assumptions specific to a given AI algorithm. Although random play does not represent expert-level gameplay, it provides a useful first approximation of the game’s depth, branching factor, balance, and dynamic behavior. These results also serve as a baseline for the agent evaluations presented later in the thesis.

Several of the metrics discussed in this chapter are inspired by previous work on game analysis, game design, and general board game concepts~\cite{shannon1950chess2}~\cite{browne2010evolutionary}~\cite{russell2010aima}~\cite{browne2008recombination}. In particular, gameplay metrics such as branching factor, game length, decisiveness, stability, and drama help characterize not only the computational difficulty of Tricolor, but also some of its structural and aesthetic properties as a game.

\section{Experimental Protocol}

In order to estimate the main properties of Tricolor, one million games were simulated between two Random Agents. In each position, a Random Agent selects uniformly among all legal moves returned by the move generator. If the forced-move rule applies, the agent therefore selects uniformly among the available capture or imprisonment moves only.

A \textit{turn} is defined as a single move made by one player. Thus, a sequence of two turns corresponds to one move by each player. This convention is used consistently throughout the thesis, including when comparing Tricolor with other games.

For each simulated game, several quantities were recorded:

\begin{itemize}
    \item the total number of turns before the game ended
    \item the number of legal actions available at each turn
    \item the final outcome of the game
    \item the evolution of position evaluations during play
    \item the number of repeated states
    \item the number of occupied board sites
    \item the average movement distance
    \item the number of pieces remaining on the board
    \item several game-dynamics metrics such as decisiveness, stability, and drama
\end{itemize}

Random simulations have two main advantages. First, they are computationally inexpensive, which makes it possible to generate a very large sample of games. Second, they do not depend on the design choices of a specific AI strategy. This makes them suitable for estimating broad structural properties of the game.

However, the results should be interpreted carefully. Random play may produce positions that strong players would avoid, and it may fail to reveal strategic patterns that only emerge under more rational play. For this reason, the metrics presented in this chapter should be understood as a baseline characterization of Tricolor, rather than as a definitive description of expert play.

\section{Computational Complexity Metrics}

\subsection{Game Duration}

The first metric measures the duration of a game in number of turns. This is important for AI agents because game length directly affects the size of the search tree and the cost of simulations. A game that tends to last many turns is generally harder to analyze exhaustively than a short game with the same branching factor.

The distribution of game durations obtained from the simulations is shown in~\ref{fig:total_turns}. The average duration is approximately 177 turns, with a standard deviation of 69 turns.

\begin{figure}[h]
    \centering
    \includegraphics[scale=0.6]{game_metrics/images/total_turns.pdf}
    \caption{Plot with the number of turns in a game}
    \label{fig:total_turns}
\end{figure}

This result shows that Tricolor games can be relatively long. For comparison, Shannon’s classical estimate for Chess uses an average game length of about 40 moves per player, corresponding to about 80 individual player moves under the convention used here~\cite{shannon1950chess2}. Under random play, Tricolor therefore appears to have a greater average depth than this classical estimate for Chess.

The high standard deviation also indicates that game duration varies substantially. Some games end quickly because captures and imprisonments rapidly reduce one player’s mobility, while others continue for much longer due to extended movement, delayed confrontation, or repeated positional reorganization.

This variability is important for AI. In Monte Carlo Tree Search, long games make random rollouts more expensive. In Alpha-Beta search, long games make it less likely that the search reaches terminal states, increasing the importance of heuristic evaluation functions.

\subsection{Branching Factor}

The branching factor measures the number of legal actions available to the player to move. It is one of the most important indicators of search complexity, since many tree-search algorithms have a cost that grows exponentially with the branching factor.

The average branching factor observed during the simulations is shown in~\ref{fig:total_branches}. Across the simulated games, Tricolor has an average branching factor of approximately 47, with a standard deviation of 9.

\begin{figure}[h]
    \centering
    \includegraphics[scale=0.6]{game_metrics/images/total_branches.pdf}
    \caption{Plot with average branching factor in a game}
    \label{fig:total_branches}
\end{figure}

This value is high. It is larger than the classical average branching factor of Chess, often estimated around 35~\cite{shannon1950chess2}. In the initial position of Tricolor, the number of legal actions is particularly high because many pieces are still available and many stacks can move in several directions. As the game progresses, pieces are captured or imprisoned, which tends to reduce the number of available actions.

The observed standard deviation reflects the diversity of game trajectories. Games in which captures occur early tend to have lower branching factors later on, because fewer active stacks remain available. Conversely, games in which players avoid major confrontations for longer tend to preserve larger numbers of legal actions.

The combination of a high branching factor and long game duration is one of the main reasons why Tricolor is computationally challenging. Even relatively shallow search depths may already involve a very large number of possible continuations.

\subsection{Game Tree Complexity}

Game tree complexity estimates the number of possible games that can be played from the initial position until termination. Computing this value exactly is infeasible for Tricolor, as it would require enumerating all legal continuations of the game. A common approximation is:

\[ GTC \approx b^d \],

where $b$ is the average branching factor and $d$ is the average game depth~\cite{shannon1950chess2}.

Using the empirical values obtained from random simulations,

\[ b\approx47 \text{ and } d\approx177, \]

we obtain:

\[ GTC \approx b^d = 47^{177} \approx 47^{1.2 \times 148} = ({47^{1.2}})^{148} \approx 100^{148}\]

Taking the base-10 logarithm gives approximately:

 \[ log_{10}(GTC) \approx 296. \]

Thus, the game tree complexity of Tricolor is roughly:

 \[ GTC \approx 10^{296} \]

This value should be interpreted as a rough estimate rather than an exact result. It depends on the average branching factor and average game length observed under random play, which may differ from values obtained under stronger play. Nevertheless, the estimate clearly indicates that Tricolor has an extremely large game tree.

This has direct consequences for AI. Exhaustive search is not realistic, and even optimized search methods can only examine a small fraction of the possible continuations. Therefore, strong play requires selective search, effective pruning, or heuristic evaluation.

\subsection{State Space Complexity}

State space complexity estimates the number of distinct positions that may occur in the game. Computing the exact number of reachable states is difficult, because not every theoretical placement of pieces corresponds to a legal position reachable from the initial state. Instead, this thesis considers an upper bound.

The board contains 61 tiles. Each player starts with 18 pieces, and pieces may be stacked on any tile. Since captured pieces are removed from the game, a player may have between 0 and 18 pieces remaining on the board.

For a given number $w$ of white pieces, the number of ways to distribute them over 61 tiles is given by the stars-and-bars formula:

 \[ \binom{60+w}{w} \]

Similarly, for $b$ black pieces, the number of possible distributions is:

 \[ \binom{60+b}{b} \]

Summing over all possible remaining numbers of white and black pieces gives:

 \[ \sum_{w=0}^{w=18}\sum_{b=0}^{b=18}\binom{60+w}{w}\binom{60+b}{b} = \sum_{w=0}^{w=18}\binom{60+w}{w} \cdot \sum_{b=0}^{b=18}\binom{60+b}{b} \]

Using the identity

 \[ \sum_{k=0}^{k=n}\binom{x+k}{k} = \binom{x+n+1}{n} \]

this becomes:

 \[ \binom{60+18+1}{18} \cdot \binom{60+18+1}{18} = \binom{79}{18}^2 \]

This distribution does not yet encode stack ownership in positions where both players have pieces on the same tile. Since at most 18 mixed stacks can exist, multiplying by $2^{18}$ gives an upper bound for ownership possibilities. Finally, the player to move adds a factor of 2. Therefore:

 \[ \binom{79}{18}^2 \cdot 2^{18} \cdot 2 = \binom{79}{18}^2 \cdot 2^{19} \]

This gives an upper bound of approximately:

 \[ UB(SSC) = \binom{79}{18}^2 \cdot 2^{19} \approx (2 \cdot 10^{15})^2 \cdot 2^{19} = 10^{30} \cdot 2^{21} \approx 10^{30} \cdot 2^{3.3 \cdot 6.4} = 10^{30} \cdot (2^{3.3})^{6.4} \approx 10^{36} \]

Therefore, the state space complexity of Tricolor is bounded above by roughly:

 \[ UB(SSC) \approx 10^{36} \]

This value is high, although still below the classical state-space estimates often given for Chess, which are around $10^{47}$~\cite{shannon1950chess2}. The result highlights an interesting distinction: Tricolor has a very large game tree because games are long and branching factors are high, but its state space remains smaller than that of some more widely studied games.

This difference is important. A large game tree makes direct search difficult, while a comparatively smaller state space suggests that future work could potentially investigate transposition tables, endgame databases, or retrograde analysis for restricted subsets of the game.

\section{Game Outcome and Balance}

The final outcome of each simulated game was recorded as a win for Player 1, a win for Player 2, or a draw. Since the initial board setup is symmetric, one important question is whether the first player has a measurable advantage under random play.

The results are shown in~\ref{fig:winrates}. Player 1 wins approximately 49.4\% of the games, while Player 2 wins approximately 48.9\%. Draws account for approximately 1.8\% of the simulated games.

\begin{figure}[h]
    \centering
    \includegraphics[scale=0.6]{game_metrics/images/winrates.pdf}
    \caption{Plot with the win percentage of each player and draws}
    \label{fig:winrates}
\end{figure}

The difference between the win rates of the two players is only 0.5 percentage points. Given the size of the sample, this suggests that there is no substantial first-player advantage under random play. This is consistent with the symmetric nature of the initial position.

The draw rate is very low. This indicates that, even when both players move randomly, most games eventually lead to a decisive outcome. The low draw rate is partly explained by the capture and imprisonment mechanics: once a player begins to lose control over stacks, their mobility may decrease rapidly, eventually leaving them with no legal moves.

However, these results should not be generalized too quickly to expert play. A low draw rate under random play does not necessarily imply a low draw rate between strong agents. Strong players may be better at preserving material, avoiding forced weaknesses, or maintaining defensive structures. The comparison with stronger AI agents is therefore revisited later in the thesis.

\section{Gameplay and Aesthetic Metrics}

In addition to classical complexity measures, this thesis also studies several gameplay-related metrics inspired by previous work on game design and general game concepts~\cite{browne2010evolutionary}~\cite{browne2008recombination}. These metrics are useful because they describe how a game evolves dynamically, rather than only how large its search space is.

To compute several of these metrics, it is necessary to define an evaluation function for a position.

\subsection{Position Evaluation}

The purpose of the position evaluation function is to estimate whether a given position favors Player 1 or Player 2. The score is normalized in the interval $[-1,1]$, where:

\begin{itemize} 
    \item 0 represents an approximately balanced position
    \item a positive score indicates an advantage for Player 1
    \item a negative score indicates an advantage for Player 2
\end{itemize}

The evaluation is based on the total controlled power of each player. Let $w_s$ be the total score of the white player and $b_s$ the total score of the black player. The position evaluation is:

\[ eval = \frac{w_s - b_s}{w_s + b_s} \]

The value $w_s$ is computed as:

\[ w_s = \sum_{i=0}^{i=60}c_i*w_i\]

where $c_i$ is the value of the color of tile $i$, and $w_i$ is the number of white pieces controlled on that tile. The value $b_s$ is computed similarly for black pieces.

This evaluation function captures two intuitive ideas. First, a player with more controlled pieces generally has more mobility and more opportunities to attack. Second, controlling pieces on stronger tiles is valuable because tile colors affect attack and defense power.

This evaluation is intentionally simple. It is not meant to represent perfect strategic understanding of Tricolor. Instead, it provides a consistent numerical basis for measuring game dynamics such as lead changes, decisiveness, and drama.

\subsection{Narrowness}

Narrowness measures how many actions are worth considering in a given position. In this thesis, an action is considered relevant if it improves or maintains the current evaluation score for the player to move.

The average narrowness observed in the simulations is shown in~\ref{fig:total_narrowness}. The mean value is approximately 25, with a standard deviation of 5.

\begin{figure}[h]
    \centering
    \includegraphics[scale=0.6]{game_metrics/images/total_narrowness.pdf}
    \caption{Plot with the average narrowness during a game}
    \label{fig:total_narrowness}
\end{figure}

Since the average branching factor is approximately 47, this means that roughly half of the available moves preserve or improve the evaluation according to the simple evaluation function. This suggests that Tricolor often presents players with a relatively large number of plausible choices.

From an AI perspective, this is important. If only a small number of moves were relevant in most positions, move ordering or pruning would be easier. Instead, Tricolor frequently offers many moves that appear reasonable according to material and positional criteria. This increases the difficulty of selecting strong actions.

The interpretation of narrowness depends strongly on the evaluation function. A stronger heuristic might classify fewer moves as genuinely relevant. Nevertheless, the result indicates that Tricolor is not a game where random positions usually contain only one or two obvious candidate moves.

\subsection{Decisiveness}

Decisiveness measures the proportion of the game completed before the eventual winner obtains an advantage that is maintained until the end. This metric is computed only for games with a winner; drawn games are excluded.

The distribution of decisiveness values is shown in~\ref{fig:decisiveness}. The average decisiveness is approximately 70\%, with a standard deviation of 22\%.

\begin{figure}[h]
    \centering
    \includegraphics[scale=0.6]{game_metrics/images/decisiveness.pdf}
    \caption{Plot with the game decisiveness}
    \label{fig:decisiveness}
\end{figure}

This relatively high value means that, on average, under random play, the eventual winner often does not establish a stable advantage until late in the game. In many games, the evaluation may fluctuate for a long time before one player finally obtains a decisive lead.

This result suggests that Tricolor has a high degree of uncertainty under random play. The outcome often remains unclear until the final phase of the game. This is consistent with the presence of imprisonment mechanics: pieces that appear lost may later be recovered if control of a stack changes.

From a gameplay perspective, this can be seen as an interesting property. Games in which the winner becomes obvious too early may feel less engaging, while games in which the evaluation remains unstable can maintain tension until the end.

\subsection{Stability}

Stability measures how often the lead changes during a game. In this thesis, a lead change occurs when the sign of the evaluation score changes from positive to negative or from negative to positive.

The percentage of lead-changing turns is shown in~\ref{fig:total_lead_change_actions}. The average value is approximately 16\%, with a standard deviation of 7\%.

\begin{figure}[h]
    \centering
    \includegraphics[scale=0.6]{game_metrics/images/total_lead_change_actions.pdf}
    \caption{Plot with the percentage of lead changing turns during a game}
    \label{fig:total_lead_change_actions}
\end{figure}

This means that in a typical random game, the apparent leader changes several times. This is consistent with the high decisiveness value reported above: if the final winner only stabilizes their advantage late in the game, then the lead is expected to fluctuate during earlier phases.

The relatively low stability of Tricolor under random play can be explained by the game’s mechanics. A single capture or imprisonment may significantly modify the evaluation, especially if it changes control of a large stack or frees imprisoned pieces. The forced-move rule may also cause sudden reversals by compelling players to enter tactical sequences.

\subsection{Drama}

Drama measures the average absolute change in evaluation between consecutive positions. It captures how abruptly the apparent advantage changes during play.

The distribution of drama values is shown in~\ref{fig:total_drama}. The average drama is approximately 0.07, with a standard deviation of 0.02.

\begin{figure}[h]
    \centering
    \includegraphics[scale=0.6]{game_metrics/images/total_drama.pdf}
    \caption{Plot with the average drama in a game}
    \label{fig:total_drama}
\end{figure}

Since the evaluation score lies between -1 and 1, an average change of 0.07 is moderate. Most individual moves do not radically transform the position. However, this does not mean that the game lacks dramatic moments. Certain tactical sequences, especially those involving imprisonment and later recovery of pieces, can produce large swings in evaluation.

Thus, the game appears to combine gradual positional evolution with occasional high-impact tactical events. This is a meaningful property for both strategic analysis and AI design: agents must evaluate not only immediate material changes, but also the possibility of future reversals caused by stack control.

\subsection{Example Game Evaluations}

To illustrate the evolution of position evaluations during play, two example games between Random Agents are shown in~\ref{fig:game_evaluation_1} and~\ref{fig:game_evaluation_2}.

\begin{figure}[h]
    \centering
    \includegraphics[scale=0.5]{game_metrics/images/game_evaluation_1.pdf}
    \caption{Plot with game 1 evaluation}
    \label{fig:game_evaluation_1}
\end{figure}

\begin{figure}[h]
    \centering
    \includegraphics[scale=0.5]{game_metrics/images/game_evaluation_2.pdf}
    \caption{Plot with game 2 evaluation}
    \label{fig:game_evaluation_2}
\end{figure}

In the first example, Player 2 eventually wins. The evaluation curve crosses the equality line several times, indicating multiple lead changes. Later, Player 1 obtains a favorable position, but Player 2 manages to recover and eventually maintain the advantage until victory. This example illustrates how Tricolor positions can remain unstable for many turns before a decisive advantage appears.

In the second example, Player 1 wins. During the first part of the game, the evaluation remains relatively close to equality. The final advantage is established only late in the game. This example is consistent with the high decisiveness values measured across the full simulation set.

These examples show that the aggregate metrics are not merely abstract quantities. They correspond to concrete gameplay patterns: fluctuating evaluations, late stabilization of advantage, and occasional dramatic shifts caused by tactical interactions.

\section{Additional Metrics}

Several additional metrics were also computed and are reported in Appendix~\ref{all_game_metrics}. These include:

\begin{itemize}
    \item Decision Moves, measuring the percentage of turns in which more than one legal move is available;
    \item State Repetition, measuring how often positions repeat during a game;
    \item Decision Factor, measuring the average number of legal actions when forced single-move turns are excluded;
    \item Board Sites Occupied, measuring how many tiles are occupied at least once during a game;
    \item Move Distance, measuring the average distance travelled by moved pieces;
    \item Piece Number, measuring the average number of pieces remaining on the board during a game;
    \item Lead Change Moves, measuring how many available actions would change the current leader.
\end{itemize}

These metrics provide additional information about the structure and dynamics of Tricolor. For example, the high percentage of decision moves indicates that most turns offer meaningful choice, while the low number of repeated states suggests that games rarely cycle extensively under random play. The board occupation and movement-distance metrics also show that games tend to use a substantial portion of the board, rather than remaining confined to the initial setup area.

The purpose of including these metrics in the appendix is to provide a broader empirical profile of the game without overloading the main analysis. They may be useful for future comparisons with other games or with alternative versions of Tricolor.

\section{Discussion}

The results presented in this chapter show that Tricolor is a computationally demanding game. Its average branching factor is high, its games are relatively long, and its estimated game tree complexity is extremely large. These properties make exhaustive search infeasible and justify the need for selective AI methods.

At the same time, the estimated state space upper bound is not as large as the game tree complexity. This suggests that Tricolor may contain many different paths leading to overlapping or related positions. Future work could investigate the use of transposition tables or state-space exploration techniques to better understand this structure.

The outcome statistics suggest that the game is relatively balanced under random play, with no clear first-player advantage. The low draw rate indicates that the capture and imprisonment mechanics usually lead to decisive outcomes, even without strategic play.

The gameplay metrics reveal an additional layer of interest. Tricolor is not only computationally large; it is also dynamically unstable. Lead changes occur frequently, decisive advantages often appear late, and tactical interactions can produce meaningful reversals. These properties suggest that the game contains rich strategic possibilities despite its limited historical study.

The analysis also highlights the limitations of random simulations. Random agents are useful for estimating baseline properties, but they do not represent strong play. Some measured values, such as game duration, branching factor, draw rate, and decisiveness, may change significantly when stronger agents are used. This is investigated later in the thesis through comparisons between Alpha-Beta and Monte Carlo Tree Search agents.

Overall, this chapter establishes Tricolor as a deep, high-branching, strategically dynamic combinatorial game. These observations motivate the next stage of the thesis: identifying strategic principles and heuristic evaluations that can guide AI agents in such a large search space.

%%%%%%%%%%%%%%%%%%%%%%%%%%%%%%%%%

\begin{comment}

This chapter presents some general metrics of Tricolor useful for better understanding the combinatorial complexity of the game. Some metrics are also a set of \textit{Aesthetic Measurements}, which have the goal of measuring how \textit{interesting} Tricolor is as a game. These concepts are defined in ~\cite{evolutionary_game_design} and ~\cite{automatic_generation_and_evaluation_of_recombination_games}.

The origin and source of inspiration of these metrics can be found in the \textit{Appendix} \ref{all_game_metrics}.

In order to reproduce the results presented in this chapter, \textbf{1 million games} have been simulated between 2 \textbf{Random Agents}. This is a convenient approach because the moves made by the agents aren't biased by specific strategies or algorithms and the games are simulated fast.

In order to provide a visual representation of the computed metrics, for most of them, a \textit{violin} plot is displayed (e.g. ~\ref{fig:total_turns}). It's a type of plot that shows the distribution of the y-values using surfaces colored in blue instead of individual dots. They are used when there are many data points that need to be represented on the graph for aesthetic reasons.

%%%%%%%%%%%%%%%%%%%%%%%%%%%%%%%%%

\section{Duration}
\label{sec:duration}

\subsection{Metric}

This metric measures the duration of a game in turns. A turn is equivalent to a move made by any of the 2 players. This metric is relevant for the performance of AI agents like \textit{Alpha-Beta} and \textit{MCTS} ~\cite{alpha_beta}. The Figure~\ref{fig:total_turns} shows the distribution of this metric measured in all the simulated games.

\begin{figure}[h]
    \centering
    \includegraphics[scale=0.6]{game_metrics/images/total_turns.pdf}
    \caption{Plot with the number of turns in a game}
    \label{fig:total_turns}
\end{figure}

\subsection{Data Interpretation}

Looking at the plot, we can make the following observations:
\begin{itemize}
    \item On average, a game lasts 177 turns. Comparing to chess, where a game lasts on average 80 turns ~\cite{shannon_chess} \footnote{In Shannon's paper about chess, a turn consists in two consecutive moves, one made by each player. However, in my measurements, a turn consists in only one move made by any player. Therefore, the number of turns as presented in Shannon's paper about chess was modified so that is satisfies my definition for the rest of this paper.}, Tricolor's games depth is much bigger. This aspect has a huge impact on the performance of AI Agents, especially MCTS, which needs to perform a high amount of random game simulations in order to find the best move.
    \item The large standard deviation ($\approx$ 69 turns) indicates that game duration varies significantly depending on the Agents’ behavior, suggesting that games do not follow a consistent flow.
\end{itemize}

%%%%%%%%%%%%%%%%%%%%%%%%%%%%%%%%%%%%%%%%

\section{Branching Factor}
\label{sec:bf}

\subsection{Metric}

This metric measures how many possible actions any player can choose from, on average, during their turn. This metric is relevant for the performance of AI agents like \textit{Alpha-Beta} and \textit{MCTS} ~\cite{alpha_beta}. The Figure~\ref{fig:total_branches} shows the distribution of this metric measured in all the simulated games.

\begin{figure}[h]
    \centering
    \includegraphics[scale=0.6]{game_metrics/images/total_branches.pdf}
    \caption{Plot with average branching factor in a game}
    \label{fig:total_branches}
\end{figure}

\subsection{Data Interpretation}

Looking at the plot, we can make the following observations:
\begin{itemize}
    \item On average, a game has a branching factor of 47. Comparing to chess, where the average branching factor is 35 ~\cite{shannon_chess}, Tricolor's branching factor is slightly higher. This aspect has a huge impact on the performance of AI Agents, especially for Alpha-Beta, as its time complexity heavily depends on the branching factor. The branching factor of the initial position of the game is 88 for both players (as it can be seen in Figure~\ref{fig:board_setup}). As the game progresses, more and more pawns are being captured and imprisoned, which reduces the number of possible actions for both players during their turns. This explains why the average branching factor is 47.
    \item The standard deviation of $\approx10$ can be explained by the standard deviation obtained previously for the games duration metric. The shorter a game is, the smaller the average branching factor is, because pieces are captured and imprisoned very early in the game. The opposite is true for longer games.
\end{itemize}

%%%%%%%%%%%%%%%%%%%%%%%%%%%%%%%%%%%%%%%%

\section{Game Complexity}

\end{comment}

\begin{comment}

\subsection{Decision Moves} % pas utile

\subsubsection{Metric}


This metric measures what percentage of the moves made by any player during a game aren't forced. A move isn't forced if the player can choose between at least 2 actions during his turn. The Figure~\ref{fig:total_more_than_one_move} shows the distribution of this metric measured in all the simulated games.

\begin{figure}[h]
    \centering
    \includegraphics[scale=0.7]{game_metrics/images/total_more_than_one_move.pdf}
    \caption{Plot with the percentage (\%) of turns where more than one move was possible during a game}
    \label{fig:total_more_than_one_move}
\end{figure}

\subsubsection{Data Interpretation}

Looking at the plot, we can make the following observations:
\begin{itemize}
    \item On average, $90\%$ of the moves made by any player during a game aren't forced. This result is expected because of one strict rule of Tricolor which says that when a player can imprison or capture an opponent stack, he has to do it.
    \item Even though this $90\%$ figure can seem quite constraining for a game (meaning the players are forced to make a move quite often), it doesn't automatically mean there are no strategies required to master it.
\end{itemize}

\end{comment}

\begin{comment}

\subsection{Game Tree Complexity}

\subsubsection{Metric}

This metric measures how many different games can be played starting from the initial position until the game ends. Two games are different if they differ by at least one action at any point. This metric is relevant for the performance of AI agents like \textit{Alpha-Beta} and \textit{MCTS} ~\cite{alpha_beta}.

Computing the exact value of this metric for Tricolor is very complex. However, a good approximation is given by the following formula ~\cite{alpha_beta}:
\[ GTC \approx b^d \],
where:
\begin{itemize}
    \item $GTC$ is the game tree complexity
    \item $b$ is the average game branching factor
    \item $d$ is the average game depth
\end{itemize}

Using the previously computed empirical values for the two game metrics in the formula, we obtain:
\[ GTC \approx b^d = 47^{177} \approx 47^{1.2 \times 144} = ({47^{1.2}})^{144} \approx 100^{144}\]

This value is astronomical. As an example, the $GTC$ of chess is $\approx 10^{120} = 100^{60}$ ~\cite{shannon_chess}, so Tricolor's $GTC$ is much larger.

\subsection{State Tree Complexity}

\subsubsection{Metric}

This metric measures how many different positions are achievable in Tricolor. In other words, it counts all the possible ways pieces can be placed on the board, including who owns each tile and which player's turn it is. This metric is relevant for the performance of AI agents like \textit{Alpha-Beta} and \textit{MCTS} ~\cite{alpha_beta}.

In order to approximate the $STC$ (Sate Tree Complexity), an upper bound is computed as it's too difficult to determine the $STC$ using a simple mathematical formula. For this, the 4 following assumptions are made so that the upper bound is as small as possible:
\begin{enumerate}
    \item Each player can have at most 18 pieces and at least 0 pieces on the board \footnote{A player can have 0 pieces if the last move consisted in a capture by the opponent.}.
    \item Each piece, no matter its color, can be placed on any of the 61 tiles.
    \item For each tile containing at least one piece of each color, 2 separate states must be considered \footnote{One for each player owning the tile}. However, counting only such tiles is too complex. In order to provide an upper bound, it is assumed each position has 18 pairs of white and black pieces on 18 different tiles. This results in $2^{18}$ different states to multiply with each previous state. 
    \item Finally, a state also encodes the player's turn. Therefore, the final number of previous states must be multiplied by 2.
\end{enumerate}

The trivial formula of combinations with repetitions \footnote{Also called \textit{stars and bars method}.} is used for the number of ways pieces can be placed on the 61 tiles. The following formula is obtained:

 \[ \text{UB}(\text{STC}) =  \underbrace{\sum_{w=0}^{w=18}\sum_{b=0}^{b=18}\binom{60+w}{w}\binom{60+b}{b}}_{\text{1st and 2nd assumptions}} \cdot  \underbrace{2^{18}}_{\text{3rd assumption}} \cdot \underbrace{2}_{\text{4th assumption}} = \]

 \[ = \sum_{w=0}^{w=18}\binom{60+w}{w} \cdot \sum_{b=0}^{b=18}\binom{60+b}{b} \cdot 2^{19} = \binom{60+18+1}{18} \cdot \binom{60+18+1}{18} \cdot 2^{19} = \]

 \[ = \binom{79}{18}^2 \cdot 2^{19} \approx (2 \cdot 10^{15})^2 \cdot 2^{19} = 10^{30} \cdot 2^{21} \approx 10^{30} \cdot 2^{3.3 \cdot 6.4} = 10^{30} \cdot (2^{3.3})^{6.4} \approx 10^{36} \] \begin{equation}
     \label{eq:STC}
 \end{equation}


Therefore, an upper bound of the $STC$ is $10^{36}$. It is notably high, but not larger than the $STC$ of chess, which is $ \approx 10^{47}$ ~\cite{shannon_chess}. 

%%%%%%%%%%%%%%%%%%%%%%%%%%%%%%%%%%%%%%%%

\end{comment}

\begin{comment}

\section{State Repetition} % utile

\subsection{Metric}

This metric measures how many times the same states were repeated in a game. In other words, a position is repeated whenever the Tricolor board has the same number of pieces on the same tiles with the same players controlling them and it's the same player's turn. The Figure~\ref{fig:total_repeated_states} shows the distribution of this metric measured in all the simulated games.

\begin{figure}[h]
    \centering
    \includegraphics[scale=0.7]{game_metrics/images/total_repeated_states.pdf}
    \caption{Plot with the number of total repeated states in a game}
    \label{fig:total_repeated_states}
\end{figure}

\subsection{Data Interpretation}

Looking at the plot, we can make the following observations:
\begin{itemize}
    \item On average, a game has 0.15 repeated states. This means that a repeated position occurs about once every 6 or 7 games, which is very rare knowing a game lasts on average 177 turns.
    \item Tricolor's mechanics are such that it's easy to repeat a position by "undoing" a move. However, the fact that the actions of the random agents lead to a new position almost every turn is explained by 2 factors: 
    \begin{enumerate}
        \item the game has a high state-space, meaning there are many achievable positions (as proved in equation~\ref{eq:STC})
        \item the branching factor is high (as shown in Figure~\ref{fig:total_branches})
    \end{enumerate}
\end{itemize}

\end{comment}

%%%%%%%%%%%%%%%%%%%%%%%%%%%%%%%%%%%%%%%%

\begin{comment}

\section{Decision Factor} % pas utile

\subsection{Metric}

This metric is very similar to the branching factor metric (see Section~\ref{sec:bf}) The only difference is when computing the average, only the turns where the players have at least 2 possible actions are considered. The Figure~\ref{fig:total_decision_factors} shows the distribution of this metric measured in all the simulated games.

\begin{figure}[h]
    \centering
    \includegraphics[scale=0.7]{game_metrics/images/total_decision_factors.pdf}
    \caption{Plot with the average decision factor during a game}
    \label{fig:total_decision_factors}
\end{figure}

\subsection{Data Interpretation}

Looking at the plot, we can make the following observations:
\begin{itemize}
    \item On average, a game has a decision factor of $ \approx 52$. This is slightly bigger than the branching factor ($\approx 47$). Moreover, the standard variation is almost the same as the one of the branching factor. This is normal and it's because turns where only one move is possible are ignored.
\end{itemize}

\end{comment}

%%%%%%%%%%%%%%%%%%%%%%%%%%%%%%%%%%%%%%%%

\begin{comment}

\section{Board Sites Occupied} % pas utile

\subsection{Metric}

This metric measures how many board tiles were occupied by any player at least once during the game. Initially, both players occupy 18 tiles each, so 36 tiles in total. The board has a total of 61 tiles, so, for each game, the value of this metric can be between 36 and 61. The Figure~\ref{fig:total_occupied_tiles} shows the distribution of this metric measured in all the simulated games.

\begin{figure}[h]
    \centering
    \includegraphics[scale=0.7]{game_metrics/images/total_occupied_tiles.pdf}
    \caption{Plot with the number of occupied tiles in a game}
    \label{fig:total_occupied_tiles}
\end{figure}

\subsection{Data Interpretation}

Looking at the plot, we can make the following observations:
\begin{itemize}
    \item The distribution has a spike shape because the number of occupied tiles is an integer. Therefore, on the graph, the spikes represent integers and the space between them is filled with some blue for a more aesthetic visual representation.
    \item The average number of occupied tiles is $\approx 54$ out of the 61, with a small standard deviation of $\approx 4$. This suggests that, for most of the games, the majority of the board tiles are used. A high board coverage is a good game feature, because it means the game is spatially active and multiple regions can matter strategically.
\end{itemize}

\end{comment}

%%%%%%%%%%%%%%%%%%%%%%%%%%%%%%%%%%%%%%%%

\begin{comment}

\section{Move Distance} % pas utile

\subsection{Metric}

This metric measures the average distance traveled by pieces when they move around the board during a game. In Tricolor, a stack of pieces can move at a distance of 1, 2 or 3 tiles, depending on the number of pieces on the stack (see Game Rules~\ref{sec:rules_movement}). The Figure~\ref{fig:total_distance} shows the distribution of this metric measured in all the simulated games.

\begin{figure}[h]
    \centering
    \includegraphics[scale=0.7]{game_metrics/images/total_distance.pdf}
    \caption{Plot with the average distance traveled by pieces in a game}
    \label{fig:total_distance}
\end{figure}

\subsection{Data Interpretation}

Looking at the plot, we can make the following observations:
\begin{itemize}
    \item The average distance traveled by stacks is $1.26$ tiles with a very small standard deviation of $0.08$. This indicates that actions consisting in moving only one piece are the most frequent. This is expected, considering that on the initial board each player occupies 18 tiles with one piece per tile.
    \item Moreover, a small average traveled distance suggests that pieces tend to be spread in smaller stacks around the board rather than condensed on a few tiles in taller stacks. This is certainly due to the high number of tiles (= 61) of the board.
\end{itemize}

\end{comment}

%%%%%%%%%%%%%%%%%%%%%%%%%%%%%%%%%%%%%%%%

\begin{comment}

\section{Piece Number} % pas utile

\subsection{Metric}

This metric measures the average number of pieces on the board during a game. Players start with 18 pieces each, so there are 36 pieces in total initially on the board. The Figure~\ref{fig:total_number_of_pieces} shows the distribution of this metric measured in all the simulated games.

\begin{figure}[h]
    \centering
    \includegraphics[scale=0.7]{game_metrics/images/total_number_of_pieces.pdf}
    \caption{Plot with the average number of pieces on the board during a game}
    \label{fig:total_number_of_pieces}
\end{figure}

\subsection{Data Interpretation}

Looking at the plot, we can make the following observations:
\begin{itemize}
    \item The average number of pieces on the board during a game is $ \approx 24$ with a standard deviation of $\approx 4$. In the previous metrics, we found out that pieces tend to be spread all around the board when considering random players. This implies that \textbf{capture} moves only remove a small amount of pieces when executed. Therefore, we can expect the number of pieces during a game to have a smooth evolution rather than an abrupt one. If the initial number of pieces on the board is 36 and the average number of pieces is 24, then we can conclude that games usually end with 12 pieces on the board (as $\frac{12 + 36}{2} = 24$).
\end{itemize}

\end{comment}

%%%%%%%%%%%%%%%%%%%%%%%%%%%%%%%%%%%%%%%%

\begin{comment}

\section{Game Outcome} % utile

\subsection{Metric}

This metric measures the percentage of games won by Player 1 (White), Player 2 (Black) and draws. The initial board setup is symmetrical, meaning no player has a positional advantage. It's useful to analyze from a \textit{game aesthetics} ~\cite{evolutionary_game_design} perspective whether making the first move brings or not an advantage in Tricolor. The Figure~\ref{fig:winrates} shows the distribution of this metric measured in all the simulated games.

\begin{figure}[h]
    \centering
    \includegraphics[scale=0.6]{game_metrics/images/winrates.pdf}
    \caption{Plot with the win percentage of each player and draws}
    \label{fig:winrates}
\end{figure}

\subsection{Data Interpretation}

Looking at the plot, we can make the following observations:
\begin{itemize}
    \item Player 1 (White) has a win percentage of $47.6\%$ and Player 2 (Black) has a win percentage of $47.3\%$, which creates a difference of $0.3\%$. Given this very small difference in win percentage, we can confidently conclude that there is no obvious starting advantage for any of the players in Tricolor.
    \item The draw rate is $\approx 7\%$, which is notably low. It indicates that even with random moves, one of the players manages to capture or imprison all of the opponent's pieces.
\end{itemize}

%%%%%%%%%%%%%%%%%%%%%%%%%%%%%%%%%%%%%%%%

\section{State Evaluation}
\label{sec:state_evaluation}

All the metrics presented in the following section have the purpose of analyzing Tricolor from a \textit{game aesthetics} perspective ~\cite{evolutionary_game_design} and ~\cite{automatic_generation_and_evaluation_of_recombination_games}.

For the following metrics, a position evaluation of the game states is needed. The goal of the evaluation is to assess whether a certain position looks equal for both players, or if one of the players has an advantage and is more likely to win from there.

For this, each position is assessed with a score taking values in the range $[-1, 1]$, where a score of 0 represents an equal position, a score $> 0$ represents an advantage for Player 1 (White) and a score $< 0$ an advantage for Player 2 (Black). Computing the score follows this formula:
\[ \frac{w_s - b_s}{w_s + b_s}, \in [-1,1]\]
The total power of the black player is subtracted from the total power of the white player and then is divided by the sum of the two (in order to normalize the score to the range $[-1, 1]$). We have:
\[ w_s = \sum_{i=0}^{i=60}c_i*w_i\]
, where $c_i$ is the number associated to the color of tile $i$ (see rules ~\ref{game_rules}) and $w_i$ is the number of controlled white pieces on tile $i$ \footnote{Imprisoned pieces are not counted.}. Similarly, $b_s$ is defined by considering the black pieces instead of the white pieces. In other words, the total power of a player is defined as the sum of the powers of all his owned stacks. The power of an owned stack is the product between the number of friendly pieces in the stack and the value associated to the color of the tile the stack sits on (see rules~\ref{sec:rules}).

This position evaluation formula is relevant because:
\begin{itemize}
    \item having more active pieces on the board results in higher chances of winning.
    \item controlling tiles with stronger colors results in more attack power and defense.
\end{itemize}

\subsection{Narrowness} % utile

\subsubsection{Metric}

This metric measures the average narrowness of all the turns of a game. The narrowness of a turn represents the number of \textit{considerable} actions a player can choose from during that turn. An action is \textit{considerable} if it improves or maintains the current score of the position for the player to move. The Figure~\ref{fig:total_narrowness} shows the distribution of this metric measured in all the simulated games.

\begin{figure}[h]
    \centering
    \includegraphics[scale=0.6]{game_metrics/images/total_narrowness.pdf}
    \caption{Plot with the average narrowness during a game}
    \label{fig:total_narrowness}
\end{figure}

\subsubsection{Data Interpretation}

Looking at the plot, we can make the following observations:
\begin{itemize}
    \item The average narrowness is $\approx 25.13$ moves with a standard deviation of $\approx 5$. Considering that the average branching factor is $\approx 47$, it means that in average about half the moves a player can choose from are worth looking at, which is a lot. This proves again that Tricolor is a complex combinatorial game.
    \item When using a strong AI Agent with a good heuristic for Tricolor (see later \ref{advanced_ai_agents}), the position evaluation is much more complex than the one used in this chapter. The purpose of this metric is mainly to show how many actions a player can consider \textit{playable} rather than truly good moves. 
\end{itemize}

\subsection{Decisiveness} % utile

\subsubsection{Metric}

This metric measures the proportion of the game completed before the eventual winner establishes a decisive advantage that is maintained until the end of the game. This metric is only computed when a game has a winning player, so draws a discarded. The Figure~\ref{fig:decisiveness} shows the distribution of this metric measured in all the simulated games.

\begin{figure}[h]
    \centering
    \includegraphics[scale=0.6]{game_metrics/images/decisiveness.pdf}
    \caption{Plot with the game decisiveness}
    \label{fig:decisiveness}
\end{figure}

\subsubsection{Data Interpretation}

Looking at the plot, we can make the following observations:
\begin{itemize}
    \item The average decisiveness is $98.82$ with a standard deviation of $6.91$. This result suggests that most games don't have an obvious winner until the end.
    \item For Random Agents players, such a high decisiveness implies that Tricolor has an \textit{uncertain} nature ~\cite{evolutionary_game_design}. A game has an \textit{uncertain} nature if the outcome of each game remains uncertain until the end, as long as both players maintain a vested interest in it, which is a good \textit{game aesthetic}.
\end{itemize}

\end{comment}

\begin{comment}

\subsection{Lead Change} % pas utile

\subsubsection{Metric}

This metric measures the average number of possible actions during a turn that change the leader of the game. The leader of the game is the player who has a positional advantage with reference to the state evaluation used. The Figure~\ref{fig:total_lead_change_moves} shows the distribution of this metric measured in all the simulated games.

\begin{figure}[h]
    \centering
    \includegraphics[scale=0.7]{game_metrics/images/total_lead_change_moves.pdf}
    \caption{Plot with the average number of lead changing moves during a game}
    \label{fig:total_lead_change_moves}
\end{figure}

\subsubsection{Data Interpretation}

Looking at the plot, we can make the following observations:
\begin{itemize}
    \item The average number of lead changing moves is $\approx 9$ with a standard deviation of $3.71$. This number is pretty high and reinforces the idea of Tricolor not having a \textbf{snowball} nature. If so many moves can change the lead, then creating a real advantage as a Random Agent player seems unlikely.
\end{itemize}

\end{comment}

\begin{comment}

\subsection{Stability} % utile

\subsubsection{Metric}

This metric measures the percentage of moves in a game that changed the leader. Its purpose is to assess how stable games are. The Figure~\ref{fig:total_lead_change_actions} shows the distribution of this metric measured in all the simulated games.

\begin{figure}[h]
    \centering
    \includegraphics[scale=0.6]{game_metrics/images/total_lead_change_actions.pdf}
    \caption{Plot with the percentage of lead changing turns during a game}
    \label{fig:total_lead_change_actions}
\end{figure}

\subsubsection{Data Interpretation}

Looking at the plot, we can make the following observations:
\begin{itemize}
    \item The average number of lead changing turns is $\approx 16\%$ with a standard deviation of $6.79$. As expected from the previous metrics, the stability of Tricolor games is low. This means that the leader tends to change frequently. It also has to do with Tricolor having an \textbf{uncertain} nature, as seen in the previous metric.
\end{itemize}

\subsection{Drama} % utile

\subsubsection{Metric}

This metric measures the average change in position evaluation scores throughout a game. The purpose is to see how drastic are the changes in consecutive positions during a game. The Figure~\ref{fig:total_drama} shows the distribution of this metric measured in all the simulated games.

\begin{figure}[h]
    \centering
    \includegraphics[scale=0.6]{game_metrics/images/total_drama.pdf}
    \caption{Plot with the average drama in a game}
    \label{fig:total_drama}
\end{figure}

\subsubsection{Data Interpretation}

Looking at the plot, we can make the following observations:
\begin{itemize}
    \item The average drama is $0.07$ with a standard deviation of $0.02$. The score of a position evaluation is in the range $[-1, 1]$, so an average change of $0.07$ is small.
    \item Even though the average game \textit{drama} is small, a player in a disastrous losing position can manage to eventually win the game. This can happen when the losing player has many imprisoned pieces in an opponent stack but later regains their control. Therefore, Tricolor games have certain positions with high \textit{drama}, which is a good game feature ~\cite{evolutionary_game_design} and ~\cite{automatic_generation_and_evaluation_of_recombination_games}.
\end{itemize}

\subsection{Game Evaluations Examples}

In this section we review and analyze two different games between Random Agents. We expect to find features that confirm the metrics found in the \textit{State Evaluation} \ref{sec:state_evaluation} section.

Each graph shows the score of each position of the game for every turn. The yellow line represents equal positions (score $\approx0$).

\subsubsection{Game 1}

In Game 1 (see Figure~\ref{fig:game_evaluation_1}), Player 1 (White) won.

Comments:
\begin{itemize}
    \item In the first 60 turns, the drama is very low. However, even though Black gets a very good position at turn 69, White manages to regain the advantage and win the game.
    \item At around turn 76, White takes the lead and maintains it until winning the game.
\end{itemize}

\begin{figure}[h]
    \centering
    \includegraphics[scale=0.5]{game_metrics/images/game_evaluation_1.pdf}
    \caption{Plot with game 1 evaluation}
    \label{fig:game_evaluation_1}
\end{figure}

\subsubsection{Game 2}

In Game 2 (see Figure~\ref{fig:game_evaluation_2}), Player 2 (Black) won.

Comments:
\begin{itemize}
    \item Throughout the game, many lead changing moves happen. This is because the yellow line intersects many times the graph of the score function.
    \item The drama starts increasing significantly at turn 140.
    \item At around turn 170, Black takes the lead and maintains it until winning the game, which happens very late in the game.
\end{itemize}

\begin{figure}[h]
    \centering
    \includegraphics[scale=0.5]{game_metrics/images/game_evaluation_2.pdf}
    \caption{Plot with game 2 evaluation}
    \label{fig:game_evaluation_2}
\end{figure}


\subsubsection{Game 3}

In Game 3 (see Figure~\ref{fig:game_evaluation_3}), it's a draw.

Comments:
\begin{itemize}
    \item Throughout the game many lead changing moves happen, especially because the game lasts many turns (around 310).
    \item The drama starts increasing and decreasing multiple times throughout the game. Looking at the graph, we can see the biggest spikes at turns 80, 180, 225 and 235.
    \item Both players had a significant advantage at multiple turns (around 100 for White and around 185 for Black), but none of them managed to make decisive moves and win the game. Therefore, the 50-moves rule was applied (see Rules~\ref{sec:rule_modification}) and the game ended in a draw.
\end{itemize}

\begin{figure}[h]
    \centering
    \includegraphics[scale=0.5]{game_metrics/images/game_evaluation_3.pdf}
    \caption{Plot with game 3 evaluation}
    \label{fig:game_evaluation_3}
\end{figure}

\end{comment}